\documentclass[11pt]{letter}
\usepackage[margin=2.5cm]{geometry}

\begin{document}

\begin{letter}{The Editors \\ \textit{The Ramanujan Journal}}

\opening{Dear Editors,}

We submit for your consideration our manuscript entitled
\textbf{``Ab initio derivation of the Berry-Keating correction coefficient
from the Conrey-Snaith triple correlation formula''}
for publication in \textit{The Ramanujan Journal}.

The gap ratio statistic of Riemann zeta zeros converges to the GUE
prediction with a correction $c/\log^2 T$, where the empirical
coefficient $c = 1.245 \pm 0.040$ was measured in a companion
paper (Paper~A). In the present work, we compute $c$ from first
principles using the Conrey-Snaith formula for the triple correlation
$R_3$, decomposing it into two channels: $c_{R_3} = -1.6$ (ab initio,
from Richardson extrapolation at $L = 1000$ to $30{,}000$) and
$c_E = +2.8$ (gap probability, by difference).

The calculation uses a parameter-free formula weighted by the
conditional gap probability $E^{\mathrm{cond}}$ (Schur complement
of the sine kernel), requiring no ad hoc threshold. We also show
that the Fredholm determinant of the Berry-Keating kernel gives the
wrong sign, establishing the kernel phase as the fundamental
obstruction that the Conrey-Snaith route bypasses.

Two companion papers provide context: the empirical measurement and
mechanism (Paper~A), and the unconditional proof that Rudnick--Sarnak
correlations and linear programming imply the Riemann Hypothesis
(Paper~C).

We confirm that this manuscript has not been published previously and is
not under consideration elsewhere.

\closing{Sincerely,}

David Alarc\'on\\
Universidad Pablo de Olavide, Sevilla, Spain\\
\texttt{dalarub@upo.es}

\end{letter}
\end{document}
